\begin{register}{H}{UART\_FIFO\_REG}{0x0000}\label{regdesc:UARTFIFOREG}
\regfield{(reserved)}{24}{8}{000000000000000000000000}%
\regfield{UART\_RXFIFO\_RD\_BYTE}{8}{0}{{0}}%
\reglabel{Reset}\regnewline%
\begin{regdesc}\begin{reglist}
\label{fielddesc:UARTRXFIFORDBYTE}\item [UART\_RXFIFO\_RD\_BYTE] UART\regindex{n} accesses FIFO via this field. (RO)
\end{reglist}\end{regdesc}
\end{register}


\begin{register}{H}{UART\_MEM\_CONF\_REG}{0x0060}\label{regdesc:UARTMEMCONFREG}
\regfield{(reserved)}{4}{28}{0000}%
\regfield{UART\_MEM\_FORCE\_PU}{1}{27}{{0}}%
\regfield{UART\_MEM\_FORCE\_PD}{1}{26}{{0}}%
\regfield{UART\_RX\_TOUT\_THRHD}{10}{16}{{0xa}}%
\regfield{UART\_RX\_FLOW\_THRHD}{9}{7}{{0x0}}%
\regfield{UART\_TX\_SIZE}{3}{4}{{0x1}}%
\regfield{UART\_RX\_SIZE}{3}{1}{{1}}%
\regfield{(reserved)}{1}{0}{0}%
\reglabel{Reset}\regnewline%
\begin{regdesc}\begin{reglist}
\label{fielddesc:UARTRXSIZE}\item [UART\_RX\_SIZE] This field is used to configure the amount of RAM allocated for RX FIFO. The default number is 128 bytes. (R/W)
\label{fielddesc:UARTTXSIZE}\item [UART\_TX\_SIZE] This field is used to configure the amount of RAM allocated for TX FIFO. The default number is 128 bytes. (R/W)
\label{fielddesc:UARTRXFLOWTHRHD}\item [UART\_RX\_FLOW\_THRHD] This field is used to configure the maximum amount of data bytes that can be received when hardware flow control works. (R/W)
\label{fielddesc:UARTRXTOUTTHRHD}\item [UART\_RX\_TOUT\_THRHD] This field is used to configure the threshold time that the receiver takes to receive one byte, in the unit of bit time (the time it takes to transfer one bit). The UART\_RXFIFO\_TOUT\_INT interrupt will be triggered when the receiver takes more time to receive one byte with UART RX\_TOUT\_EN set to 1. (R/W)
\label{fielddesc:UARTMEMFORCEPD}\item [UART\_MEM\_FORCE\_PD] Set this bit to force power down UART RAM. (R/W)
\label{fielddesc:UARTMEMFORCEPU}\item [UART\_MEM\_FORCE\_PU] Set this bit to force power up UART RAM. (R/W)
\end{reglist}\end{regdesc}
\end{register}


\begin{register}{H}{UART\_INT\_RAW\_REG}{0x0004}\label{regdesc:UARTINTRAWREG}
\regfield{(reserved)}{12}{20}{000000000000}%
\regfield{UART\_WAKEUP\_INT\_RAW}{1}{19}{{0}}%
\regfield{UART\_AT\_CMD\_CHAR\_DET\_INT\_RAW}{1}{18}{{0}}%
\regfield{UART\_RS485\_CLASH\_INT\_RAW}{1}{17}{{0}}%
\regfield{UART\_RS485\_FRM\_ERR\_INT\_RAW}{1}{16}{{0}}%
\regfield{UART\_RS485\_PARITY\_ERR\_INT\_RAW}{1}{15}{{0}}%
\regfield{UART\_TX\_DONE\_INT\_RAW}{1}{14}{{0}}%
\regfield{UART\_TX\_BRK\_IDLE\_DONE\_INT\_RAW}{1}{13}{{0}}%
\regfield{UART\_TX\_BRK\_DONE\_INT\_RAW}{1}{12}{{0}}%
\regfield{UART\_GLITCH\_DET\_INT\_RAW}{1}{11}{{0}}%
\regfield{UART\_SW\_XOFF\_INT\_RAW}{1}{10}{{0}}%
\regfield{UART\_SW\_XON\_INT\_RAW}{1}{9}{{0}}%
\regfield{UART\_RXFIFO\_TOUT\_INT\_RAW}{1}{8}{{0}}%
\regfield{UART\_BRK\_DET\_INT\_RAW}{1}{7}{{0}}%
\regfield{UART\_CTS\_CHG\_INT\_RAW}{1}{6}{{0}}%
\regfield{UART\_DSR\_CHG\_INT\_RAW}{1}{5}{{0}}%
\regfield{UART\_RXFIFO\_OVF\_INT\_RAW}{1}{4}{{0}}%
\regfield{UART\_FRM\_ERR\_INT\_RAW}{1}{3}{{0}}%
\regfield{UART\_PARITY\_ERR\_INT\_RAW}{1}{2}{{0}}%
\regfield{UART\_TXFIFO\_EMPTY\_INT\_RAW}{1}{1}{{1}}%
\regfield{UART\_RXFIFO\_FULL\_INT\_RAW}{1}{0}{{0}}%
\reglabel{Reset}\regnewline%
\begin{regdesc}\begin{reglist}
\label{fielddesc:UARTRXFIFOFULLINTRAW}\item [UART\_RXFIFO\_FULL\_INT\_RAW] This interrupt raw bit turns to high level when the receiver receives more data than what UART\_RXFIFO\_FULL\_THRHD specifies. (R/WTC/SS)
\label{fielddesc:UARTTXFIFOEMPTYINTRAW}\item [UART\_TXFIFO\_EMPTY\_INT\_RAW] This interrupt raw bit turns to high level when the amount of data in TX FIFO is less than what UART\_TXFIFO\_EMPTY\_THRHD specifies. (R/WTC/SS)
\label{fielddesc:UARTPARITYERRINTRAW}\item [UART\_PARITY\_ERR\_INT\_RAW] This interrupt raw bit turns to high level when the receiver detects a parity error in the data. (R/WTC/SS)
\label{fielddesc:UARTFRMERRINTRAW}\item [UART\_FRM\_ERR\_INT\_RAW] This interrupt raw bit turns to high level when the receiver detects a data frame error. (R/WTC/SS)
\label{fielddesc:UARTRXFIFOOVFINTRAW}\item [UART\_RXFIFO\_OVF\_INT\_RAW] This interrupt raw bit turns to high level when the receiver receives more data than the capacity of RX FIFO. (R/WTC/SS)
\label{fielddesc:UARTDSRCHGINTRAW}\item [UART\_DSR\_CHG\_INT\_RAW] This interrupt raw bit turns to high level when the receiver detects the edge change of DSRn signal. (R/WTC/SS)
\label{fielddesc:UARTCTSCHGINTRAW}\item [UART\_CTS\_CHG\_INT\_RAW] This interrupt raw bit turns to high level when the receiver detects the edge change of CTSn signal. (R/WTC/SS)
\label{fielddesc:UARTBRKDETINTRAW}\item [UART\_BRK\_DET\_INT\_RAW] This interrupt raw bit turns to high level when the receiver detects a 0 after the stop bit. (R/WTC/SS)
\label{fielddesc:UARTRXFIFOTOUTINTRAW}\item [UART\_RXFIFO\_TOUT\_INT\_RAW] This interrupt raw bit turns to high level when the receiver takes more time than UART\_RX\_TOUT\_THRHD to receive a byte. (R/WTC/SS)
\label{fielddesc:UARTSWXONINTRAW}\item [UART\_SW\_XON\_INT\_RAW] This interrupt raw bit turns to high level when the receiver receives an XON character and UART\_SW\_FLOW\_CON\_EN is set to 1. (R/WTC/SS)
\label{fielddesc:UARTSWXOFFINTRAW}\item [UART\_SW\_XOFF\_INT\_RAW] This interrupt raw bit turns to high level when the receiver receives an XOFF character and UART\_SW\_FLOW\_CON\_EN is set to 1. (R/WTC/SS)
\label{fielddesc:UARTGLITCHDETINTRAW}\item [UART\_GLITCH\_DET\_INT\_RAW] This interrupt raw bit turns to high level when the receiver detects a glitch in the middle of a start bit. (R/WTC/SS)

\item [Continued on the next page...]
\end{reglist}\end{regdesc}
\end{register}

\addtocounter{Regfloat}{-1}
\begin{register}{H}{UART\_INT\_RAW\_REG}{0x0004}
\begin{regdesc}\begin{reglist}
\item [Continued from the previous page...]

\label{fielddesc:UARTTXBRKDONEINTRAW}\item [UART\_TX\_BRK\_DONE\_INT\_RAW] This interrupt raw bit turns to high level when the transmitter completes  sending  NULL characters, after all data in TX FIFO are sent. (R/WTC/SS)
\label{fielddesc:UARTTXBRKIDLEDONEINTRAW}\item [UART\_TX\_BRK\_IDLE\_DONE\_INT\_RAW] This interrupt raw bit turns to high level when the transmitter has kept the shortest duration after sending the last data. (R/WTC/SS)
\label{fielddesc:UARTTXDONEINTRAW}\item [UART\_TX\_DONE\_INT\_RAW] This interrupt raw bit turns to high level when the transmitter has sent out all data in FIFO. (R/WTC/SS)
\label{fielddesc:UARTRS485PARITYERRINTRAW}\item [UART\_RS485\_PARITY\_ERR\_INT\_RAW] This interrupt raw bit turns to high level when the receiver detects a parity error from the echo of the transmitter in RS485 mode. (R/WTC/SS)
\label{fielddesc:UARTRS485FRMERRINTRAW}\item [UART\_RS485\_FRM\_ERR\_INT\_RAW] This interrupt raw bit turns to high level when the receiver detects a data frame error from the echo of the transmitter in RS485 mode. (R/WTC/SS)
\label{fielddesc:UARTRS485CLASHINTRAW}\item [UART\_RS485\_CLASH\_INT\_RAW] This interrupt raw bit turns to high level when a collision is detected between the transmitter and the receiver in RS485 mode. (R/WTC/SS)
\label{fielddesc:UARTATCMDCHARDETINTRAW}\item [UART\_AT\_CMD\_CHAR\_DET\_INT\_RAW] This interrupt raw bit turns to high level when the receiver detects the configured UART\_AT\_CMD\_CHAR. (R/WTC/SS)
\label{fielddesc:UARTWAKEUPINTRAW}\item [UART\_WAKEUP\_INT\_RAW] This interrupt raw bit turns to high level when the input RXD edge changes more times than what (UART\_ACTIVE\_THRESHOLD + 3) specifies in Light-sleep mode. (R/WTC/SS)
\end{reglist}\end{regdesc}
\end{register}


\begin{register}{H}{UART\_INT\_ST\_REG}{0x0008}\label{regdesc:UARTINTSTREG}
\regfield{(reserved)}{12}{20}{000000000000}%
\regfield{UART\_WAKEUP\_INT\_ST}{1}{19}{{0}}%
\regfield{UART\_AT\_CMD\_CHAR\_DET\_INT\_ST}{1}{18}{{0}}%
\regfield{UART\_RS485\_CLASH\_INT\_ST}{1}{17}{{0}}%
\regfield{UART\_RS485\_FRM\_ERR\_INT\_ST}{1}{16}{{0}}%
\regfield{UART\_RS485\_PARITY\_ERR\_INT\_ST}{1}{15}{{0}}%
\regfield{UART\_TX\_DONE\_INT\_ST}{1}{14}{{0}}%
\regfield{UART\_TX\_BRK\_IDLE\_DONE\_INT\_ST}{1}{13}{{0}}%
\regfield{UART\_TX\_BRK\_DONE\_INT\_ST}{1}{12}{{0}}%
\regfield{UART\_GLITCH\_DET\_INT\_ST}{1}{11}{{0}}%
\regfield{UART\_SW\_XOFF\_INT\_ST}{1}{10}{{0}}%
\regfield{UART\_SW\_XON\_INT\_ST}{1}{9}{{0}}%
\regfield{UART\_RXFIFO\_TOUT\_INT\_ST}{1}{8}{{0}}%
\regfield{UART\_BRK\_DET\_INT\_ST}{1}{7}{{0}}%
\regfield{UART\_CTS\_CHG\_INT\_ST}{1}{6}{{0}}%
\regfield{UART\_DSR\_CHG\_INT\_ST}{1}{5}{{0}}%
\regfield{UART\_RXFIFO\_OVF\_INT\_ST}{1}{4}{{0}}%
\regfield{UART\_FRM\_ERR\_INT\_ST}{1}{3}{{0}}%
\regfield{UART\_PARITY\_ERR\_INT\_ST}{1}{2}{{0}}%
\regfield{UART\_TXFIFO\_EMPTY\_INT\_ST}{1}{1}{{0}}%
\regfield{UART\_RXFIFO\_FULL\_INT\_ST}{1}{0}{{0}}%
\reglabel{Reset}\regnewline%
\begin{regdesc}\begin{reglist}
\label{fielddesc:UARTRXFIFOFULLINTST}\item [UART\_RXFIFO\_FULL\_INT\_ST] This is the status bit for UART\_RXFIFO\_FULL\_INT when UART\_RXFIFO\_FULL\_INT\_ENA is set to 1. (RO)
\label{fielddesc:UARTTXFIFOEMPTYINTST}\item [UART\_TXFIFO\_EMPTY\_INT\_ST] This is the status bit for UART\_TXFIFO\_EMPTY\_INT  when UART\_TXFIFO\_EMPTY\_INT\_ENA is set to 1. (RO)
\label{fielddesc:UARTPARITYERRINTST}\item [UART\_PARITY\_ERR\_INT\_ST] This is the status bit for UART\_PARITY\_ERR\_INT when UART\_PARITY\_ERR\_INT\_ENA is set to 1. (RO)
\label{fielddesc:UARTFRMERRINTST}\item [UART\_FRM\_ERR\_INT\_ST] This is the status bit for UART\_FRM\_ERR\_INT when UART\_FRM\_ERR\_INT\_ENA is set to 1. (RO)
\label{fielddesc:UARTRXFIFOOVFINTST}\item [UART\_RXFIFO\_OVF\_INT\_ST] This is the status bit for UART\_RXFIFO\_OVF\_INT when UART\_RXFIFO\_OVF\_INT\_ENA is set to 1. (RO)
\label{fielddesc:UARTDSRCHGINTST}\item [UART\_DSR\_CHG\_INT\_ST] This is the status bit for UART\_DSR\_CHG\_INT when UART\_DSR\_CHG\_INT\_ENA is set to 1. (RO)
\label{fielddesc:UARTCTSCHGINTST}\item [UART\_CTS\_CHG\_INT\_ST] This is the status bit for UART\_CTS\_CHG\_INT when UART\_CTS\_CHG\_INT\_ENA is set to 1. (RO)
\label{fielddesc:UARTBRKDETINTST}\item [UART\_BRK\_DET\_INT\_ST] This is the status bit for UART\_BRK\_DET\_INT when UART\_BRK\_DET\_INT\_ENA is set to 1. (RO)
\label{fielddesc:UARTRXFIFOTOUTINTST}\item [UART\_RXFIFO\_TOUT\_INT\_ST] This is the status bit for UART\_RXFIFO\_TOUT\_INT when UART\_RXFIFO\_TOUT\_INT\_ENA is set to 1. (RO)
\label{fielddesc:UARTSWXONINTST}\item [UART\_SW\_XON\_INT\_ST] This is the status bit for UART\_SW\_XON\_INT when UART\_SW\_XON\_INT\_ENA is set to 1. (RO)
\label{fielddesc:UARTSWXOFFINTST}\item [UART\_SW\_XOFF\_INT\_ST] This is the status bit for UART\_SW\_XOFF\_INT when UART\_SW\_XOFF\_INT\_ENA is set to 1. (RO)
\label{fielddesc:UARTGLITCHDETINTST}\item [UART\_GLITCH\_DET\_INT\_ST] This is the status bit for UART\_GLITCH\_DET\_INT when UART\_GLITCH\_DET\_INT\_ENA is set to 1. (RO)
\label{fielddesc:UARTTXBRKDONEINTST}\item [UART\_TX\_BRK\_DONE\_INT\_ST] This is the status bit for UART\_TX\_BRK\_DONE\_INT when UART\_TX\_BRK\_DONE\_INT\_ENA is set to 1. (RO)

\item [Continued on the next page...]
\end{reglist}\end{regdesc}
\end{register}

\addtocounter{Regfloat}{-1}
\begin{register}{H}{UART\_INT\_ST\_REG}{0x0008}
\begin{regdesc}\begin{reglist}
\item [Continued from the previous page...]

\label{fielddesc:UARTTXBRKIDLEDONEINTST}\item [UART\_TX\_BRK\_IDLE\_DONE\_INT\_ST] This is the status bit for UART\_TX\_BRK\_IDLE\_DONE\_INT when UART\_TX\_BRK\_IDLE\_DONE\_INT\_ENA is set to 1. (RO)
\label{fielddesc:UARTTXDONEINTST}\item [UART\_TX\_DONE\_INT\_ST] This is the status bit for UART\_TX\_DONE\_INT when UART\_TX\_DONE\_INT\_ENA is set to 1. (RO)
\label{fielddesc:UARTRS485PARITYERRINTST}\item [UART\_RS485\_PARITY\_ERR\_INT\_ST] This is the status bit for UART\_RS485\_PARITY\_ERR\_INT when UART\_RS485\_PARITY\_INT\_ENA is set to 1. (RO)
\label{fielddesc:UARTRS485FRMERRINTST}\item [UART\_RS485\_FRM\_ERR\_INT\_ST] This is the status bit for UART\_RS485\_FRM\_ERR\_INT when UART\_RS485\_FRM\_ERR\_INT\_ENA is set to 1. (RO)
\label{fielddesc:UARTRS485CLASHINTST}\item [UART\_RS485\_CLASH\_INT\_ST] This is the status bit for UART\_RS485\_CLASH\_INT when UART\_RS485\_CLASH\_INT\_ENA is set to 1. (RO)
\label{fielddesc:UARTATCMDCHARDETINTST}\item [UART\_AT\_CMD\_CHAR\_DET\_INT\_ST] This is the status bit for UART\_AT\_CMD\_CHAR\_DET\_INT when UART\_AT\_CMD\_CHAR\_DET\_INT\_ENA is set to 1. (RO)
\label{fielddesc:UARTWAKEUPINTST}\item [UART\_WAKEUP\_INT\_ST] This is the status bit for UART\_WAKEUP\_INT when UART\_WAKEUP\_INT\_ENA is set to 1. (RO)
\end{reglist}\end{regdesc}
\end{register}


\begin{register}{H}{UART\_INT\_ENA\_REG}{0x000C}\label{regdesc:UARTINTENAREG}
\regfield{(reserved)}{12}{20}{000000000000}%
\regfield{UART\_WAKEUP\_INT\_ENA}{1}{19}{{0}}%
\regfield{UART\_AT\_CMD\_CHAR\_DET\_INT\_ENA}{1}{18}{{0}}%
\regfield{UART\_RS485\_CLASH\_INT\_ENA}{1}{17}{{0}}%
\regfield{UART\_RS485\_FRM\_ERR\_INT\_ENA}{1}{16}{{0}}%
\regfield{UART\_RS485\_PARITY\_ERR\_INT\_ENA}{1}{15}{{0}}%
\regfield{UART\_TX\_DONE\_INT\_ENA}{1}{14}{{0}}%
\regfield{UART\_TX\_BRK\_IDLE\_DONE\_INT\_ENA}{1}{13}{{0}}%
\regfield{UART\_TX\_BRK\_DONE\_INT\_ENA}{1}{12}{{0}}%
\regfield{UART\_GLITCH\_DET\_INT\_ENA}{1}{11}{{0}}%
\regfield{UART\_SW\_XOFF\_INT\_ENA}{1}{10}{{0}}%
\regfield{UART\_SW\_XON\_INT\_ENA}{1}{9}{{0}}%
\regfield{UART\_RXFIFO\_TOUT\_INT\_ENA}{1}{8}{{0}}%
\regfield{UART\_BRK\_DET\_INT\_ENA}{1}{7}{{0}}%
\regfield{UART\_CTS\_CHG\_INT\_ENA}{1}{6}{{0}}%
\regfield{UART\_DSR\_CHG\_INT\_ENA}{1}{5}{{0}}%
\regfield{UART\_RXFIFO\_OVF\_INT\_ENA}{1}{4}{{0}}%
\regfield{UART\_FRM\_ERR\_INT\_ENA}{1}{3}{{0}}%
\regfield{UART\_PARITY\_ERR\_INT\_ENA}{1}{2}{{0}}%
\regfield{UART\_TXFIFO\_EMPTY\_INT\_ENA}{1}{1}{{0}}%
\regfield{UART\_RXFIFO\_FULL\_INT\_ENA}{1}{0}{{0}}%
\reglabel{Reset}\regnewline%
\begin{regdesc}\begin{reglist}
\label{fielddesc:UARTRXFIFOFULLINTENA}\item [UART\_RXFIFO\_FULL\_INT\_ENA] This is the enable bit for UART\_RXFIFO\_FULL\_INT. (R/W)
\label{fielddesc:UARTTXFIFOEMPTYINTENA}\item [UART\_TXFIFO\_EMPTY\_INT\_ENA] This is the enable bit for UART\_TXFIFO\_EMPTY\_INT. (R/W)
\label{fielddesc:UARTPARITYERRINTENA}\item [UART\_PARITY\_ERR\_INT\_ENA] This is the enable bit for UART\_PARITY\_ERR\_INT. (R/W)
\label{fielddesc:UARTFRMERRINTENA}\item [UART\_FRM\_ERR\_INT\_ENA] This is the enable bit for UART\_FRM\_ERR\_INT. (R/W)
\label{fielddesc:UARTRXFIFOOVFINTENA}\item [UART\_RXFIFO\_OVF\_INT\_ENA] This is the enable bit for UART\_RXFIFO\_OVF\_INT. (R/W)
\label{fielddesc:UARTDSRCHGINTENA}\item [UART\_DSR\_CHG\_INT\_ENA] This is the enable bit for UART\_DSR\_CHG\_INT. (R/W)
\label{fielddesc:UARTCTSCHGINTENA}\item [UART\_CTS\_CHG\_INT\_ENA] This is the enable bit for UART\_CTS\_CHG\_INT. (R/W)
\label{fielddesc:UARTBRKDETINTENA}\item [UART\_BRK\_DET\_INT\_ENA] This is the enable bit for UART\_BRK\_DET\_INT. (R/W)
\label{fielddesc:UARTRXFIFOTOUTINTENA}\item [UART\_RXFIFO\_TOUT\_INT\_ENA] This is the enable bit for UART\_RXFIFO\_TOUT\_INT. (R/W)
\label{fielddesc:UARTSWXONINTENA}\item [UART\_SW\_XON\_INT\_ENA] This is the enable bit for UART\_SW\_XON\_INT. (R/W)
\label{fielddesc:UARTSWXOFFINTENA}\item [UART\_SW\_XOFF\_INT\_ENA] This is the enable bit for UART\_SW\_XOFF\_INT. (R/W)
\label{fielddesc:UARTGLITCHDETINTENA}\item [UART\_GLITCH\_DET\_INT\_ENA] This is the enable bit for UART\_GLITCH\_DET\_INT. (R/W)
\label{fielddesc:UARTTXBRKDONEINTENA}\item [UART\_TX\_BRK\_DONE\_INT\_ENA] This is the enable bit for UART\_TX\_BRK\_DONE\_INT. (R/W)
\label{fielddesc:UARTTXBRKIDLEDONEINTENA}\item [UART\_TX\_BRK\_IDLE\_DONE\_INT\_ENA] This is the enable bit for UART\_TX\_BRK\_IDLE\_DONE\_INT. (R/W)
\label{fielddesc:UARTTXDONEINTENA}\item [UART\_TX\_DONE\_INT\_ENA] This is the enable bit for UART\_TX\_DONE\_INT. (R/W)

\item [Continued on the next page...]
\end{reglist}\end{regdesc}
\end{register}

\addtocounter{Regfloat}{-1}
\begin{register}{H}{UART\_INT\_ENA\_REG}{0x000C}
\begin{regdesc}\begin{reglist}
\item [Continued from the previous page...]


\label{fielddesc:UARTRS485PARITYERRINTENA}\item [UART\_RS485\_PARITY\_ERR\_INT\_ENA] This is the enable bit for UART\_RS485\_PARITY\_ERR\_INT. (R/W)
\label{fielddesc:UARTRS485FRMERRINTENA}\item [UART\_RS485\_FRM\_ERR\_INT\_ENA] This is the enable bit for UART\_RS485\_PARITY\_ERR\_INT. (R/W)
\label{fielddesc:UARTRS485CLASHINTENA}\item [UART\_RS485\_CLASH\_INT\_ENA] This is the enable bit for UART\_RS485\_CLASH\_INT. (R/W)
\label{fielddesc:UARTATCMDCHARDETINTENA}\item [UART\_AT\_CMD\_CHAR\_DET\_INT\_ENA] This is the enable bit for UART\_AT\_CMD\_CHAR\_DET\_INT. (R/W)
\label{fielddesc:UARTWAKEUPINTENA}\item [UART\_WAKEUP\_INT\_ENA] This is the enable bit for UART\_WAKEUP\_INT. (R/W)
\end{reglist}\end{regdesc}
\end{register}


\begin{register}{H}{UART\_INT\_CLR\_REG}{0x0010}\label{regdesc:UARTINTCLRREG}
\regfield{(reserved)}{12}{20}{000000000000}%
\regfield{UART\_WAKEUP\_INT\_CLR}{1}{19}{{0}}%
\regfield{UART\_AT\_CMD\_CHAR\_DET\_INT\_CLR}{1}{18}{{0}}%
\regfield{UART\_RS485\_CLASH\_INT\_CLR}{1}{17}{{0}}%
\regfield{UART\_RS485\_FRM\_ERR\_INT\_CLR}{1}{16}{{0}}%
\regfield{UART\_RS485\_PARITY\_ERR\_INT\_CLR}{1}{15}{{0}}%
\regfield{UART\_TX\_DONE\_INT\_CLR}{1}{14}{{0}}%
\regfield{UART\_TX\_BRK\_IDLE\_DONE\_INT\_CLR}{1}{13}{{0}}%
\regfield{UART\_TX\_BRK\_DONE\_INT\_CLR}{1}{12}{{0}}%
\regfield{UART\_GLITCH\_DET\_INT\_CLR}{1}{11}{{0}}%
\regfield{UART\_SW\_XOFF\_INT\_CLR}{1}{10}{{0}}%
\regfield{UART\_SW\_XON\_INT\_CLR}{1}{9}{{0}}%
\regfield{UART\_RXFIFO\_TOUT\_INT\_CLR}{1}{8}{{0}}%
\regfield{UART\_BRK\_DET\_INT\_CLR}{1}{7}{{0}}%
\regfield{UART\_CTS\_CHG\_INT\_CLR}{1}{6}{{0}}%
\regfield{UART\_DSR\_CHG\_INT\_CLR}{1}{5}{{0}}%
\regfield{UART\_RXFIFO\_OVF\_INT\_CLR}{1}{4}{{0}}%
\regfield{UART\_FRM\_ERR\_INT\_CLR}{1}{3}{{0}}%
\regfield{UART\_PARITY\_ERR\_INT\_CLR}{1}{2}{{0}}%
\regfield{UART\_TXFIFO\_EMPTY\_INT\_CLR}{1}{1}{{0}}%
\regfield{UART\_RXFIFO\_FULL\_INT\_CLR}{1}{0}{{0}}%
\reglabel{Reset}\regnewline%
\begin{regdesc}\begin{reglist}
\label{fielddesc:UARTRXFIFOFULLINTCLR}\item [UART\_RXFIFO\_FULL\_INT\_CLR] Set this bit to clear the UART\_THE RXFIFO\_FULL\_INT interrupt. (WT)
\label{fielddesc:UARTTXFIFOEMPTYINTCLR}\item [UART\_TXFIFO\_EMPTY\_INT\_CLR] Set this bit to clear the UART\_TXFIFO\_EMPTY\_INT interrupt. (WT)
\label{fielddesc:UARTPARITYERRINTCLR}\item [UART\_PARITY\_ERR\_INT\_CLR] Set this bit to clear the UART\_PARITY\_ERR\_INT interrupt. (WT)
\label{fielddesc:UARTFRMERRINTCLR}\item [UART\_FRM\_ERR\_INT\_CLR] Set this bit to clear the UART\_FRM\_ERR\_INT interrupt. (WT)
\label{fielddesc:UARTRXFIFOOVFINTCLR}\item [UART\_RXFIFO\_OVF\_INT\_CLR] Set this bit to clear the UART\_UART\_RXFIFO\_OVF\_INT interrupt. (WT)
\label{fielddesc:UARTDSRCHGINTCLR}\item [UART\_DSR\_CHG\_INT\_CLR] Set this bit to clear the UART\_DSR\_CHG\_INT interrupt. (WT)
\label{fielddesc:UARTCTSCHGINTCLR}\item [UART\_CTS\_CHG\_INT\_CLR] Set this bit to clear the UART\_CTS\_CHG\_INT interrupt. (WT)
\label{fielddesc:UARTBRKDETINTCLR}\item [UART\_BRK\_DET\_INT\_CLR] Set this bit to clear the UART\_BRK\_DET\_INT interrupt. (WT)
\label{fielddesc:UARTRXFIFOTOUTINTCLR}\item [UART\_RXFIFO\_TOUT\_INT\_CLR] Set this bit to clear the UART\_RXFIFO\_TOUT\_INT interrupt. (WT)
\label{fielddesc:UARTSWXONINTCLR}\item [UART\_SW\_XON\_INT\_CLR] Set this bit to clear the UART\_SW\_XON\_INT interrupt. (WT)
\label{fielddesc:UARTSWXOFFINTCLR}\item [UART\_SW\_XOFF\_INT\_CLR] Set this bit to clear the UART\_SW\_XOFF\_INT interrupt. (WT)
\label{fielddesc:UARTGLITCHDETINTCLR}\item [UART\_GLITCH\_DET\_INT\_CLR] Set this bit to clear the UART\_GLITCH\_DET\_INT interrupt. (WT)
\label{fielddesc:UARTTXBRKDONEINTCLR}\item [UART\_TX\_BRK\_DONE\_INT\_CLR] Set this bit to clear the UART\_TX\_BRK\_DONE\_INT interrupt. (WT)
\label{fielddesc:UARTTXBRKIDLEDONEINTCLR}\item [UART\_TX\_BRK\_IDLE\_DONE\_INT\_CLR] Set this bit to clear the UART\_TX\_BRK\_IDLE\_DONE\_INT interrupt. (WT)
\label{fielddesc:UARTTXDONEINTCLR}\item [UART\_TX\_DONE\_INT\_CLR] Set this bit to clear the UART\_TX\_DONE\_INT interrupt. (WT)
\label{fielddesc:UARTRS485PARITYERRINTCLR}\item [UART\_RS485\_PARITY\_ERR\_INT\_CLR] Set this bit to clear the UART\_RS485\_PARITY\_ERR\_INT interrupt. (WT)

\item [Continued on the next page...]
\end{reglist}\end{regdesc}
\end{register}

\addtocounter{Regfloat}{-1}
\begin{register}{H}{UART\_INT\_CLR\_REG}{0x0010}
\begin{regdesc}\begin{reglist}
\item [Continued from the previous page...]

\label{fielddesc:UARTRS485FRMERRINTCLR}\item [UART\_RS485\_FRM\_ERR\_INT\_CLR] Set this bit to clear the UART\_RS485\_FRM\_ERR\_INT interrupt. (WT)
\label{fielddesc:UARTRS485CLASHINTCLR}\item [UART\_RS485\_CLASH\_INT\_CLR] Set this bit to clear the UART\_RS485\_CLASH\_INT interrupt. (WT)
\label{fielddesc:UARTATCMDCHARDETINTCLR}\item [UART\_AT\_CMD\_CHAR\_DET\_INT\_CLR] Set this bit to clear the UART\_AT\_CMD\_CHAR\_DET\_INT interrupt. (WT)
\label{fielddesc:UARTWAKEUPINTCLR}\item [UART\_WAKEUP\_INT\_CLR] Set this bit to clear the UART\_WAKEUP\_INT interrupt. (WT)
\end{reglist}\end{regdesc}
\end{register}


\begin{register}{H}{UART\_CLKDIV\_REG}{0x0014}\label{regdesc:UARTCLKDIVREG}
\regfield{(reserved)}{8}{24}{00000000}%
\regfield{UART\_CLKDIV\_FRAG}{4}{20}{{0x0}}%
\regfield{(reserved)}{8}{12}{00000000}%
\regfield{UART\_CLKDIV}{12}{0}{{0x2b6}}%
\reglabel{Reset}\regnewline%
\begin{regdesc}\begin{reglist}
\label{fielddesc:UARTCLKDIV}\item [UART\_CLKDIV] The integral part of the frequency divisor. (R/W)
\label{fielddesc:UARTCLKDIVFRAG}\item [UART\_CLKDIV\_FRAG] The fractional part of the frequency divisor. (R/W)
\end{reglist}\end{regdesc}
\end{register}


\begin{register}{H}{UART\_RX\_FILT\_REG}{0x0018}\label{regdesc:UARTRXFILTREG}
\regfield{(reserved)}{23}{9}{00000000000000000000000}%
\regfield{UART\_GLITCH\_FILT\_EN}{1}{8}{{0}}%
\regfield{UART\_GLITCH\_FILT}{8}{0}{{0x8}}%
\reglabel{Reset}\regnewline%
\begin{regdesc}\begin{reglist}
\label{fielddesc:UARTGLITCHFILT}\item [UART\_GLITCH\_FILT] When input pulse width is lower than this value, the pulse is ignored. (R/W)
\label{fielddesc:UARTGLITCHFILTEN}\item [UART\_GLITCH\_FILT\_EN] Set this bit to enable RX signal filter. (R/W)
\end{reglist}\end{regdesc}
\end{register}


\begin{register}{H}{UART\_CONF0\_REG}{0x0020}\label{regdesc:UARTCONF0REG}
\regfield{(reserved)}{3}{29}{000}%
\regfield{UART\_MEM\_CLK\_EN}{1}{28}{{1}}%
\regfield{UART\_AUTOBAUD\_EN}{1}{27}{{0}}%
\regfield{UART\_ERR\_WR\_MASK}{1}{26}{{0}}%
\regfield{UART\_CLK\_EN}{1}{25}{{0}}%
\regfield{UART\_DTR\_INV}{1}{24}{{0}}%
\regfield{UART\_RTS\_INV}{1}{23}{{0}}%
\regfield{UART\_TXD\_INV}{1}{22}{{0}}%
\regfield{UART\_DSR\_INV}{1}{21}{{0}}%
\regfield{UART\_CTS\_INV}{1}{20}{{0}}%
\regfield{UART\_RXD\_INV}{1}{19}{{0}}%
\regfield{UART\_TXFIFO\_RST}{1}{18}{{0}}%
\regfield{UART\_RXFIFO\_RST}{1}{17}{{0}}%
\regfield{UART\_IRDA\_EN}{1}{16}{{0}}%
\regfield{UART\_TX\_FLOW\_EN}{1}{15}{{0}}%
\regfield{UART\_LOOPBACK}{1}{14}{{0}}%
\regfield{UART\_IRDA\_RX\_INV}{1}{13}{{0}}%
\regfield{UART\_IRDA\_TX\_INV}{1}{12}{{0}}%
\regfield{UART\_IRDA\_WCTL}{1}{11}{{0}}%
\regfield{UART\_IRDA\_TX\_EN}{1}{10}{{0}}%
\regfield{UART\_IRDA\_DPLX}{1}{9}{{0}}%
\regfield{UART\_TXD\_BRK}{1}{8}{{0}}%
\regfield{UART\_SW\_DTR}{1}{7}{{0}}%
\regfield{UART\_SW\_RTS}{1}{6}{{0}}%
\regfield{UART\_STOP\_BIT\_NUM}{2}{4}{{1}}%
\regfield{UART\_BIT\_NUM}{2}{2}{{3}}%
\regfield{UART\_PARITY\_EN}{1}{1}{{0}}%
\regfield{UART\_PARITY}{1}{0}{{0}}%
\reglabel{Reset}\regnewline%
\begin{regdesc}\begin{reglist}
\label{fielddesc:UARTPARITY}\item [UART\_PARITY] This bit is used to configure the parity check mode. (R/W)
\label{fielddesc:UARTPARITYEN}\item [UART\_PARITY\_EN] Set this bit to enable UART parity check. (R/W)
\label{fielddesc:UARTBITNUM}\item [UART\_BIT\_NUM] This field is used to set the length of data. (R/W)
\label{fielddesc:UARTSTOPBITNUM}\item [UART\_STOP\_BIT\_NUM] This field is used to set the length of  stop bit. (R/W)
\label{fielddesc:UARTSWRTS}\item [UART\_SW\_RTS] This bit is used to configure the software RTS signal which is used in software flow control. (R/W)
\label{fielddesc:UARTSWDTR}\item [UART\_SW\_DTR] This bit is used to configure the software DTR signal which is used in software flow control. (R/W)
\label{fielddesc:UARTTXDBRK}\item [UART\_TXD\_BRK] Set this bit to enable the transmitter to send NULL characters when the process of sending data is done. (R/W)
\label{fielddesc:UARTIRDADPLX}\item [UART\_IRDA\_DPLX] Set this bit to enable IrDA loopback mode. (R/W)
\label{fielddesc:UARTIRDATXEN}\item [UART\_IRDA\_TX\_EN] This is the start enable bit for IrDA transmitter. (R/W)
\label{fielddesc:UARTIRDAWCTL}\item [UART\_IRDA\_WCTL] 1: The IrDA transmitter's 11th bit is the same as 10th bit; 0: Set IrDA transmitter's 11th bit to 0. (R/W)
\label{fielddesc:UARTIRDATXINV}\item [UART\_IRDA\_TX\_INV] Set this bit to invert the level of IrDA transmitter. (R/W)
\label{fielddesc:UARTIRDARXINV}\item [UART\_IRDA\_RX\_INV] Set this bit to invert the level of IrDA receiver. (R/W)
\label{fielddesc:UARTLOOPBACK}\item [UART\_LOOPBACK] Set this bit to enable UART loopback test mode. (R/W)
\label{fielddesc:UARTTXFLOWEN}\item [UART\_TX\_FLOW\_EN] Set this bit to enable flow control function for the transmitter. (R/W)
\label{fielddesc:UARTIRDAEN}\item [UART\_IRDA\_EN] Set this bit to enable IrDA protocol. (R/W)
\label{fielddesc:UARTRXFIFORST}\item [UART\_RXFIFO\_RST] Set this bit to reset the UART RX FIFO. (R/W)
\label{fielddesc:UARTTXFIFORST}\item [UART\_TXFIFO\_RST] Set this bit to reset the UART TX FIFO. (R/W)
\label{fielddesc:UARTRXDINV}\item [UART\_RXD\_INV] Set this bit to invert the level value of UART RXD signal. (R/W)
\label{fielddesc:UARTCTSINV}\item [UART\_CTS\_INV] Set this bit to invert the level value of UART CTS signal. (R/W)
\label{fielddesc:UARTDSRINV}\item [UART\_DSR\_INV] Set this bit to invert the level value of UART DSR signal. (R/W)


\item [Continued on the next page...]
\end{reglist}\end{regdesc}
\end{register}

\addtocounter{Regfloat}{-1}
\begin{register}{H}{UART\_CONF0\_REG}{0x0020}
\begin{regdesc}\begin{reglist}
\item [Continued from the previous page...]


\label{fielddesc:UARTTXDINV}\item [UART\_TXD\_INV] Set this bit to invert the level value of UART TXD signal. (R/W)
\label{fielddesc:UARTRTSINV}\item [UART\_RTS\_INV] Set this bit to invert the level value of UART RTS signal. (R/W)
\label{fielddesc:UARTDTRINV}\item [UART\_DTR\_INV] Set this bit to invert the level value of UART DTR signal. (R/W)
\label{fielddesc:UARTCLKEN}\item [UART\_CLK\_EN] 1: Force clock on for register; 0: Support clock only when application writes registers. (R/W)
\label{fielddesc:UARTERRWRMASK}\item [UART\_ERR\_WR\_MASK] 1: The receiver stops storing data into FIFO when data is wrong; 0: The receiver stores the data even if the  received data is wrong. (R/W)
\label{fielddesc:UARTAUTOBAUDEN}\item [UART\_AUTOBAUD\_EN] This is the enable bit for baud rate detection. (R/W)
\label{fielddesc:UARTMEMCLKEN}\item [UART\_MEM\_CLK\_EN] The signal to enable UART RAM clock gating. (R/W)
\end{reglist}\end{regdesc}
\end{register}


\begin{register}{H}{UART\_CONF1\_REG}{0x0024}\label{regdesc:UARTCONF1REG}
\regfield{(reserved)}{10}{22}{0000000000}%
\regfield{UART\_RX\_TOUT\_EN}{1}{21}{{0}}%
\regfield{UART\_RX\_FLOW\_EN}{1}{20}{{0}}%
\regfield{UART\_RX\_TOUT\_FLOW\_DIS}{1}{19}{{0}}%
\regfield{UART\_DIS\_RX\_DAT\_OVF}{1}{18}{{0}}%
\regfield{UART\_TXFIFO\_EMPTY\_THRHD}{9}{9}{{0x60}}%
\regfield{UART\_RXFIFO\_FULL\_THRHD}{9}{0}{{0x60}}%
\reglabel{Reset}\regnewline%
\begin{regdesc}\begin{reglist}
\label{fielddesc:UARTRXFIFOFULLTHRHD}\item [UART\_RXFIFO\_FULL\_THRHD] An UART\_RXFIFO\_FULL\_INT interrupt is generated when the receiver receives more data than the value of this field. (R/W)
\label{fielddesc:UARTTXFIFOEMPTYTHRHD}\item [UART\_TXFIFO\_EMPTY\_THRHD] An UART\_TXFIFO\_EMPTY\_INT interrupt is generated when the number of data bytes in TX FIFO is less than the value of this field. (R/W)
\label{fielddesc:UARTDISRXDATOVF}\item [UART\_DIS\_RX\_DAT\_OVF] Disable UART RX data overflow detection. (R/W)
\label{fielddesc:UARTRXTOUTFLOWDIS}\item [UART\_RX\_TOUT\_FLOW\_DIS] Set this bit to stop accumulating idle\_cnt when hardware flow control works. (R/W)
\label{fielddesc:UARTRXFLOWEN}\item [UART\_RX\_FLOW\_EN] This is the flow enable bit for UART receiver. (R/W)
\label{fielddesc:UARTRXTOUTEN}\item [UART\_RX\_TOUT\_EN] This is the enable bit for UART receiver's timeout function. (R/W)
\end{reglist}\end{regdesc}
\end{register}


\begin{register}{H}{UART\_FLOW\_CONF\_REG}{0x0034}\label{regdesc:UARTFLOWCONFREG}
\regfield{(reserved)}{26}{6}{00000000000000000000000000}%
\regfield{UART\_SEND\_XOFF}{1}{5}{{0}}%
\regfield{UART\_SEND\_XON}{1}{4}{{0}}%
\regfield{UART\_FORCE\_XOFF}{1}{3}{{0}}%
\regfield{UART\_FORCE\_XON}{1}{2}{{0}}%
\regfield{UART\_XONOFF\_DEL}{1}{1}{{0}}%
\regfield{UART\_SW\_FLOW\_CON\_EN}{1}{0}{{0}}%
\reglabel{Reset}\regnewline%
\begin{regdesc}\begin{reglist}
\label{fielddesc:UARTSWFLOWCONEN}\item [UART\_SW\_FLOW\_CON\_EN] Set this bit to enable software flow control. When UART receives flow control characters XON or XOFF, which can be configured by UART\_XON\_CHAR or UART\_XOFF\_CHAR respectively, UART\_SW\_XON\_INT or UART\_SW\_XOFF\_INT interrupts can be triggered if enabled.  (R/W)
\label{fielddesc:UARTXONOFFDEL}\item [UART\_XONOFF\_DEL] Set this bit to remove flow control characters from the received data. (R/W)
\label{fielddesc:UARTFORCEXON}\item [UART\_FORCE\_XON] Set this bit to force the transmitter to send data. (R/W)
\label{fielddesc:UARTFORCEXOFF}\item [UART\_FORCE\_XOFF] Set this bit to stop the transmitter from sending data. (R/W)
\label{fielddesc:UARTSENDXON}\item [UART\_SEND\_XON] Set this bit to send an XON character. This bit is cleared by hardware automatically. (R/W/SS/SC)
\label{fielddesc:UARTSENDXOFF}\item [UART\_SEND\_XOFF] Set this bit to send an XOFF character. This bit is cleared by hardware automatically. (R/W/SS/SC)
\end{reglist}\end{regdesc}
\end{register}


\begin{register}{H}{UART\_SLEEP\_CONF\_REG}{0x0038}\label{regdesc:UARTSLEEPCONFREG}
\regfield{(reserved)}{22}{10}{0000000000000000000000}%
\regfield{UART\_ACTIVE\_THRESHOLD}{10}{0}{{0xf0}}%
\reglabel{Reset}\regnewline%
\begin{regdesc}\begin{reglist}
\label{fielddesc:UARTACTIVETHRESHOLD}\item [UART\_ACTIVE\_THRESHOLD] UART is activated from Light-sleep mode when the input RXD edge changes more times than the value of this field plus 3. (R/W)
\end{reglist}\end{regdesc}
\end{register}


\begin{register}{H}{UART\_SWFC\_CONF0\_REG}{0x003C}\label{regdesc:UARTSWFCCONF0REG}
\regfield{(reserved)}{15}{17}{000000000000000}%
\regfield{UART\_XOFF\_CHAR}{8}{9}{{0x13}}%
\regfield{UART\_XOFF\_THRESHOLD}{9}{0}{{0xe0}}%
\reglabel{Reset}\regnewline%
\begin{regdesc}\begin{reglist}
\label{fielddesc:UARTXOFFTHRESHOLD}\item [UART\_XOFF\_THRESHOLD] When the number of data bytes in RX FIFO is more than the value of this field with UART\_SW\_FLOW\_CON\_EN set to 1, the transmitter sends an XOFF character. (R/W)
\label{fielddesc:UARTXOFFCHAR}\item [UART\_XOFF\_CHAR] This field stores the XOFF flow control character. (R/W)
\end{reglist}\end{regdesc}
\end{register}


\begin{register}{H}{UART\_SWFC\_CONF1\_REG}{0x0040}\label{regdesc:UARTSWFCCONF1REG}
\regfield{(reserved)}{15}{17}{000000000000000}%
\regfield{UART\_XON\_CHAR}{8}{9}{{0x11}}%
\regfield{UART\_XON\_THRESHOLD}{9}{0}{{0x0}}%
\reglabel{Reset}\regnewline%
\begin{regdesc}\begin{reglist}
\label{fielddesc:UARTXONTHRESHOLD}\item [UART\_XON\_THRESHOLD] When the number of data bytes in RX FIFO is less than the value of this field with UART\_SW\_FLOW\_CON\_EN set to 1, the transmitter sends an XON character. (R/W)
\label{fielddesc:UARTXONCHAR}\item [UART\_XON\_CHAR] This field stores the XON flow control character. (R/W)
\end{reglist}\end{regdesc}
\end{register}


\begin{register}{H}{UART\_TXBRK\_CONF\_REG}{0x0044}\label{regdesc:UARTTXBRKCONFREG}
\regfield{(reserved)}{24}{8}{000000000000000000000000}%
\regfield{UART\_TX\_BRK\_NUM}{8}{0}{{0xa}}%
\reglabel{Reset}\regnewline%
\begin{regdesc}\begin{reglist}
\label{fielddesc:UARTTXBRKNUM}\item [UART\_TX\_BRK\_NUM] This field is used to configure the number of 0 to be sent after the process of sending data is done. It is active when UART\_TXD\_BRK is set to 1. (R/W)
\end{reglist}\end{regdesc}
\end{register}


\begin{register}{H}{UART\_IDLE\_CONF\_REG}{0x0048}\label{regdesc:UARTIDLECONFREG}
\regfield{(reserved)}{12}{20}{000000000000}%
\regfield{UART\_TX\_IDLE\_NUM}{10}{10}{{0x100}}%
\regfield{UART\_RX\_IDLE\_THRHD}{10}{0}{{0x100}}%
\reglabel{Reset}\regnewline%
\begin{regdesc}\begin{reglist}
\label{fielddesc:UARTRXIDLETHRHD}\item [UART\_RX\_IDLE\_THRHD] A frame end signal is generated when the receiver takes more time to receive one byte data than the value of this field, in the unit of bit time (the time it takes to transfer one bit). (R/W)
\label{fielddesc:UARTTXIDLENUM}\item [UART\_TX\_IDLE\_NUM] This field is used to configure the duration time between transfers, in the unit of bit time (the time it takes to transfer one bit). (R/W)
\end{reglist}\end{regdesc}
\end{register}


\begin{register}{H}{UART\_RS485\_CONF\_REG}{0x004C}\label{regdesc:UARTRS485CONFREG}
\regfield{(reserved)}{22}{10}{0000000000000000000000}%
\regfield{UART\_RS485\_TX\_DLY\_NUM}{4}{6}{{0}}%
\regfield{UART\_RS485\_RX\_DLY\_NUM}{1}{5}{{0}}%
\regfield{UART\_RS485RXBY\_TX\_EN}{1}{4}{{0}}%
\regfield{UART\_RS485TX\_RX\_EN}{1}{3}{{0}}%
\regfield{UART\_DL1\_EN}{1}{2}{{0}}%
\regfield{UART\_DL0\_EN}{1}{1}{{0}}%
\regfield{UART\_RS485\_EN}{1}{0}{{0}}%
\reglabel{Reset}\regnewline%
\begin{regdesc}\begin{reglist}
\label{fielddesc:UARTRS485EN}\item [UART\_RS485\_EN] Set this bit to choose RS485 mode. (R/W)
\label{fielddesc:UARTDL0EN}\item [UART\_DL0\_EN] Configures whether or not to add a turnaround delay of 1 bit before the start bit.\\
0: Not add\\
1: Add\\(R/W)
\label{fielddesc:UARTDL1EN}\item [UART\_DL1\_EN] Configures whether or not to add a turnaround delay of 1 bit after the stop bit.\\
0: Not add\\
1: Add\\(R/W)
\label{fielddesc:UARTRS485TXRXEN}\item [UART\_RS485TX\_RX\_EN] Set this bit to enable the receiver could receive data when the transmitter is transmitting data in RS485 mode. (R/W)
\label{fielddesc:UARTRS485RXBYTXEN}\item [UART\_RS485RXBY\_TX\_EN] 1: enable RS485 transmitter to send data when RS485 receiver line is busy. (R/W)
\label{fielddesc:UARTRS485RXDLYNUM}\item [UART\_RS485\_RX\_DLY\_NUM] This bit is used to delay the receiver's internal data signal. (R/W)
\label{fielddesc:UARTRS485TXDLYNUM}\item [UART\_RS485\_TX\_DLY\_NUM] This field is used to delay the transmitter's internal data signal. (R/W)
\end{reglist}\end{regdesc}
\end{register}


\begin{register}{H}{UART\_CLK\_CONF\_REG}{0x0078}\label{regdesc:UARTCLKCONFREG}
\regfield{(reserved)}{6}{26}{000000}%
\regfield{UART\_RX\_SCLK\_EN}{1}{25}{{1}}%
\regfield{UART\_TX\_SCLK\_EN}{1}{24}{{1}}%
\regfield{UART\_RST\_CORE}{1}{23}{{0}}%
\regfield{UART\_SCLK\_EN}{1}{22}{{1}}%
\regfield{UART\_SCLK\_SEL}{2}{20}{{3}}%
\regfield{UART\_SCLK\_DIV\_NUM}{8}{12}{{0x1}}%
\regfield{UART\_SCLK\_DIV\_A}{6}{6}{{0x0}}%
\regfield{UART\_SCLK\_DIV\_B}{6}{0}{{0x0}}%
\reglabel{Reset}\regnewline%
\begin{regdesc}\begin{reglist}
\label{fielddesc:UARTSCLKDIVB}\item [UART\_SCLK\_DIV\_B] The denominator of the frequency divisor. (R/W)
\label{fielddesc:UARTSCLKDIVA}\item [UART\_SCLK\_DIV\_A] The numerator of the frequency divisor. (R/W)
\label{fielddesc:UARTSCLKDIVNUM}\item [UART\_SCLK\_DIV\_NUM] The integral part of the frequency divisor. (R/W)
\label{fielddesc:UARTSCLKSEL}\item [UART\_SCLK\_SEL] Selects UART clock source. 1: APB\_CLK; 2: RC\_FAST\_CLK; 3: XTAL\_CLK. (R/W)
\label{fielddesc:UARTSCLKEN}\item [UART\_SCLK\_EN] Set this bit to enable UART TX/RX clock. (R/W)
\label{fielddesc:UARTRSTCORE}\item [UART\_RST\_CORE] Write 1 and then write 0 to this bit, to reset UART TX/RX. (R/W)
\label{fielddesc:UARTTXSCLKEN}\item [UART\_TX\_SCLK\_EN] Set this bit to enable UART TX clock. (R/W)
\label{fielddesc:UARTRXSCLKEN}\item [UART\_RX\_SCLK\_EN] Set this bit to enable UART RX clock. (R/W)
\end{reglist}\end{regdesc}
\end{register}


\begin{register}{H}{UART\_STATUS\_REG}{0x001C}\label{regdesc:UARTSTATUSREG}
\regfield{UART\_TXD}{1}{31}{{1}}%
\regfield{UART\_RTSN}{1}{30}{{1}}%
\regfield{UART\_DTRN}{1}{29}{{1}}%
\regfield{(reserved)}{3}{26}{000}%
\regfield{UART\_TXFIFO\_CNT}{10}{16}{{0}}%
\regfield{UART\_RXD}{1}{15}{{1}}%
\regfield{UART\_CTSN}{1}{14}{{1}}%
\regfield{UART\_DSRN}{1}{13}{{0}}%
\regfield{(reserved)}{3}{10}{000}%
\regfield{UART\_RXFIFO\_CNT}{10}{0}{{0}}%
\reglabel{Reset}\regnewline%
\begin{regdesc}\begin{reglist}
\label{fielddesc:UARTRXFIFOCNT}\item [UART\_RXFIFO\_CNT] Stores the number of valid data bytes in RX FIFO. (RO)
\label{fielddesc:UARTDSRN}\item [UART\_DSRN] This bit represents the level of the internal UART DSR signal. (RO)
\label{fielddesc:UARTCTSN}\item [UART\_CTSN] This bit represents the level of the internal UART CTS signal. (RO)
\label{fielddesc:UARTRXD}\item [UART\_RXD] This bit represents the level of the internal UART RXD signal. (RO)
\label{fielddesc:UARTTXFIFOCNT}\item [UART\_TXFIFO\_CNT] Stores the number of data bytes in TX FIFO. (RO)
\label{fielddesc:UARTDTRN}\item [UART\_DTRN] This bit represents the level of the internal UART DTR signal. (RO)
\label{fielddesc:UARTRTSN}\item [UART\_RTSN] This bit represents the level of the internal UART RTS signal. (RO)
\label{fielddesc:UARTTXD}\item [UART\_TXD] This bit represents the  level of the internal UART TXD signal. (RO)
\end{reglist}\end{regdesc}
\end{register}


\begin{register}{H}{UART\_MEM\_TX\_STATUS\_REG}{0x0064}\label{regdesc:UARTMEMTXSTATUSREG}
\regfield{(reserved)}{11}{21}{00000000000}%
\regfield{UART\_TX\_RADDR}{10}{11}{{0x0}}%
\regfield{(reserved)}{1}{10}{0}%
\regfield{UART\_APB\_TX\_WADDR}{10}{0}{{0x0}}%
\reglabel{Reset}\regnewline%
\begin{regdesc}\begin{reglist}
\label{fielddesc:UARTAPBTXWADDR}\item [UART\_APB\_TX\_WADDR] This field stores the offset address in TX FIFO when software writes TX FIFO via APB. (RO)
\label{fielddesc:UARTTXRADDR}\item [UART\_TX\_RADDR] This field stores the offset address in TX FIFO when TX FSM reads data via Tx\_FIFO\_Ctrl. (RO)
\end{reglist}\end{regdesc}
\end{register}


\begin{register}{H}{UART\_MEM\_RX\_STATUS\_REG}{0x0068}\label{regdesc:UARTMEMRXSTATUSREG}
\regfield{(reserved)}{11}{21}{00000000000}%
\regfield{UART\_RX\_WADDR}{10}{11}{{0x100}}%
\regfield{(reserved)}{1}{10}{0}%
\regfield{UART\_APB\_RX\_RADDR}{10}{0}{{0x100}}%
\reglabel{Reset}\regnewline%
\begin{regdesc}\begin{reglist}
\label{fielddesc:UARTAPBRXRADDR}\item [UART\_APB\_RX\_RADDR] This field stores the offset address in RX FIFO when software reads data from RX FIFO via APB. UART0 is 0x200. UART1 is 0x280. (RO)
\label{fielddesc:UARTRXWADDR}\item [UART\_RX\_WADDR] This field stores the offset address in RX FIFO when Rx\_FIFO\_Ctrl writes RX FIFO. (RO)
\end{reglist}\end{regdesc}
\end{register}


\begin{register}{H}{UART\_FSM\_STATUS\_REG}{0x006C}\label{regdesc:UARTFSMSTATUSREG}
\regfield{(reserved)}{24}{8}{000000000000000000000000}%
\regfield{UART\_ST\_UTX\_OUT}{4}{4}{{0}}%
\regfield{UART\_ST\_URX\_OUT}{4}{0}{{0}}%
\reglabel{Reset}\regnewline%
\begin{regdesc}\begin{reglist}
\label{fielddesc:UARTSTURXOUT}\item [UART\_ST\_URX\_OUT] This is the status field of the receiver. (RO)
\label{fielddesc:UARTSTUTXOUT}\item [UART\_ST\_UTX\_OUT] This is the status field of the transmitter. (RO)
\end{reglist}\end{regdesc}
\end{register}


\begin{register}{H}{UART\_LOWPULSE\_REG}{0x0028}\label{regdesc:UARTLOWPULSEREG}
\regfield{(reserved)}{20}{12}{00000000000000000000}%
\regfield{UART\_LOWPULSE\_MIN\_CNT}{12}{0}{{0xfff}}%
\reglabel{Reset}\regnewline%
\begin{regdesc}\begin{reglist}
\label{fielddesc:UARTLOWPULSEMINCNT}\item [UART\_LOWPULSE\_MIN\_CNT] This field stores the value of the minimum duration time of the low level pulse, in the unit of APB\_CLK cycles. It is used in baud rate detection. (RO)
\end{reglist}\end{regdesc}
\end{register}


\begin{register}{H}{UART\_HIGHPULSE\_REG}{0x002C}\label{regdesc:UARTHIGHPULSEREG}
\regfield{(reserved)}{20}{12}{00000000000000000000}%
\regfield{UART\_HIGHPULSE\_MIN\_CNT}{12}{0}{{0xfff}}%
\reglabel{Reset}\regnewline%
\begin{regdesc}\begin{reglist}
\label{fielddesc:UARTHIGHPULSEMINCNT}\item [UART\_HIGHPULSE\_MIN\_CNT] This field stores the value of the maximum duration time for the high level pulse, in the unit of APB\_CLK cycles. It is used in baud rate detection. (RO)
\end{reglist}\end{regdesc}
\end{register}


\begin{register}{H}{UART\_RXD\_CNT\_REG}{0x0030}\label{regdesc:UARTRXDCNTREG}
\regfield{(reserved)}{22}{10}{0000000000000000000000}%
\regfield{UART\_RXD\_EDGE\_CNT}{10}{0}{{0x0}}%
\reglabel{Reset}\regnewline%
\begin{regdesc}\begin{reglist}
\label{fielddesc:UARTRXDEDGECNT}\item [UART\_RXD\_EDGE\_CNT] This field stores the count of RXD edge change.  It is used in baud rate detection. (RO)
\end{reglist}\end{regdesc}
\end{register}


\begin{register}{H}{UART\_POSPULSE\_REG}{0x0070}\label{regdesc:UARTPOSPULSEREG}
\regfield{(reserved)}{20}{12}{00000000000000000000}%
\regfield{UART\_POSEDGE\_MIN\_CNT}{12}{0}{{0xfff}}%
\reglabel{Reset}\regnewline%
\begin{regdesc}\begin{reglist}
\label{fielddesc:UARTPOSEDGEMINCNT}\item [UART\_POSEDGE\_MIN\_CNT] This field stores the minimal input clock count between two positive edges. It is used in baud rate detection. (RO)
\end{reglist}\end{regdesc}
\end{register}


\begin{register}{H}{UART\_NEGPULSE\_REG}{0x0074}\label{regdesc:UARTNEGPULSEREG}
\regfield{(reserved)}{20}{12}{00000000000000000000}%
\regfield{UART\_NEGEDGE\_MIN\_CNT}{12}{0}{{0xfff}}%
\reglabel{Reset}\regnewline%
\begin{regdesc}\begin{reglist}
\label{fielddesc:UARTNEGEDGEMINCNT}\item [UART\_NEGEDGE\_MIN\_CNT] This field stores the minimal input clock count between two negative edges. It is used in baud rate detection. (RO)
\end{reglist}\end{regdesc}
\end{register}


\begin{register}{H}{UART\_AT\_CMD\_PRECNT\_REG}{0x0050}\label{regdesc:UARTATCMDPRECNTREG}
\regfield{(reserved)}{16}{16}{0000000000000000}%
\regfield{UART\_PRE\_IDLE\_NUM}{16}{0}{{0x901}}%
\reglabel{Reset}\regnewline%
\begin{regdesc}\begin{reglist}
\label{fielddesc:UARTPREIDLENUM}\item [UART\_PRE\_IDLE\_NUM] This field is used to configure the idle duration time before the first AT\_CMD is received by the receiver, in the unit of bit time (the time it takes to transfer one bit). (R/W)
\end{reglist}\end{regdesc}
\end{register}


\begin{register}{H}{UART\_AT\_CMD\_POSTCNT\_REG}{0x0054}\label{regdesc:UARTATCMDPOSTCNTREG}
\regfield{(reserved)}{16}{16}{0000000000000000}%
\regfield{UART\_POST\_IDLE\_NUM}{16}{0}{{0x901}}%
\reglabel{Reset}\regnewline%
\begin{regdesc}\begin{reglist}
\label{fielddesc:UARTPOSTIDLENUM}\item [UART\_POST\_IDLE\_NUM] This field is used to configure the duration time between the last AT\_CMD and the next data byte, in the unit of bit time (the time it takes to transfer one bit). (R/W)
\end{reglist}\end{regdesc}
\end{register}


\begin{register}{H}{UART\_AT\_CMD\_GAPTOUT\_REG}{0x0058}\label{regdesc:UARTATCMDGAPTOUTREG}
\regfield{(reserved)}{16}{16}{0000000000000000}%
\regfield{UART\_RX\_GAP\_TOUT}{16}{0}{{11}}%
\reglabel{Reset}\regnewline%
\begin{regdesc}\begin{reglist}
\label{fielddesc:UARTRXGAPTOUT}\item [UART\_RX\_GAP\_TOUT] This field is used to configure the duration time between the AT\_CMD characters, in the unit of bit time (the time it takes to transfer one bit). (R/W)
\end{reglist}\end{regdesc}
\end{register}


\begin{register}{H}{UART\_AT\_CMD\_CHAR\_REG}{0x005C}\label{regdesc:UARTATCMDCHARREG}
\regfield{(reserved)}{16}{16}{0000000000000000}%
\regfield{UART\_CHAR\_NUM}{8}{8}{{0x3}}%
\regfield{UART\_AT\_CMD\_CHAR}{8}{0}{{0x2b}}%
\reglabel{Reset}\regnewline%
\begin{regdesc}\begin{reglist}
\label{fielddesc:UARTATCMDCHAR}\item [UART\_AT\_CMD\_CHAR] This field is used to configure the content of AT\_CMD character. (R/W)
\label{fielddesc:UARTCHARNUM}\item [UART\_CHAR\_NUM] This field is used to configure the number of continuous AT\_CMD characters received by the receiver. (R/W)
\end{reglist}\end{regdesc}
\end{register}


\begin{register}{H}{UART\_DATE\_REG}{0x007C}\label{regdesc:UARTDATEREG}
\regfield{UART\_DATE}{32}{0}{{0x2008270}}%
\reglabel{Reset}\regnewline%
\begin{regdesc}\begin{reglist}
\label{fielddesc:UARTDATE}\item [UART\_DATE] This is the version control register. (R/W)
\end{reglist}\end{regdesc}
\end{register}


\begin{register}{H}{UART\_ID\_REG}{0x0080}\label{regdesc:UARTIDREG}
\regfield{UART\_REG\_UPDATE}{1}{31}{{0}}%
\regfield{UART\_UPDATE\_CTRL}{1}{30}{{1}}%
\regfield{UART\_ID}{30}{0}{{0x000500}}%
\reglabel{Reset}\regnewline%
\begin{regdesc}\begin{reglist}
\label{fielddesc:UARTID}\item [UART\_ID] This field is used to configure the UART\_ID. (R/W)
\label{fielddesc:UARTUPDATECTRL}\item [UART\_UPDATE\_CTRL] This bit is used to control register synchronization mode. This bit must be cleared before writing 1 to UART\_REG\_UPDATE to synchronize configured values to UART Core's clock domain. (R/W)
\label{fielddesc:UARTREGUPDATE}\item [UART\_REG\_UPDATE] When this bit is set to 1 by software, registers are synchronized to UART Core's clock domain. This bit is cleared by hardware after synchronization is done. (R/W/SC)
\end{reglist}\end{regdesc}
\end{register}
